% theory_classic_methods.tex —— 潮流计算 通用理论：经典方法（GS/Z-bus 谱系、FDPF 解耦、直流潮流、DistFlow 与前推回代、分布式松弛与自适应孤岛）
% 本文件为理论层章节，遵循 docs/modules/THEORY_WRITING_GUIDE.md。

\section{经典潮流方法理论：不动点迭代、解耦与线性化}\label{sec:pf-th-cls}

\begin{theorynote}
本章为通用数学理论，按经典教材与原始文献体系撰写，覆盖五组主题：
(i) Gauss--Seidel 与阻抗阵法的历史定位及其作为不动点迭代的数学结构；
(ii) 快速解耦潮流（FDPF）的解耦物理原理推导、XB/BX 变体差异与收敛性分析；
(iii) 直流潮流从 AC 方程经三条假设到 $P=B\theta$ 的严格推导、线性化误差界、损耗近似，
及基于直流模型的转移分布因子（PTDF/ISF）与线路开断分布因子（LODF）；
(iv) 辐射状配电网的 DistFlow 支路潮流模型、前推回代的不动点--压缩映射收敛理论
与高 $R/X$ 弱网特性；
(v) 分布式松弛（含方程内精确加边 Newton 系统）与自适应孤岛的理论定位；
(vi) 以平台求解器实测数据给出的 FDPF、直流与前推回代收敛性数值基准。
Newton 法本体、全局化策略与电压稳定/换流器相关理论各有专章，本章不复述；
本章公式不要求逐式对应代码。与平台实现（\srcpath{src/power\_flow/}）的对应关系
见本章末"与实现的对应"表，超出实现的内容以"未实现扩展说明"框就地标记。
\end{theorynote}

\subsection{记号与约定}\label{sec:pf-th-cls-notation}
\begin{itemize}
  \item $n$：AC 母线数；母线复电压 $V_i=V_{m,i}\,e^{j\theta_i}$（极坐标），
        $\theta_{ij}=\theta_i-\theta_j$；
  \item $Y_{bus}=G+jB$：节点导纳阵，$Y_{ij}=G_{ij}+jB_{ij}$；支路 $k$（$i\!\to\!j$）
        阻抗 $z_k=r_k+jx_k$，串联导纳 $y_k=g_k+jb_k=1/z_k$（感应性支路 $b_k<0$）；
  \item $P_i,Q_i$：母线 $i$ 注入有功/无功（发电减负荷，pu）；上标
        $\mathrm{spec}/\mathrm{calc}$ 区分指定值与计算值；
        失配 $\Delta P_i=P_i^{\mathrm{spec}}-P_i^{\mathrm{calc}}$，
        $\Delta Q_i$ 同理；
  \item $B',B''\in\mathbb R^{n\times n}$：FDPF 的两个常量近似阵
        （本章取\textbf{ Laplacian 口径}：对角元为正，迭代方程写作
        $B'\Delta\theta=\Delta P/V$、$B''\Delta V=\Delta Q/V$；
        部分教材采用反号口径 $-B'$，两者仅差整体符号约定）；
  \item $\rho(\cdot)$：谱半径；$\lVert\cdot\rVert_\infty$：行和范数（向量时为最大模）；
  \item 配网一节中，$\mathcal T$ 为以根（slack）为源的辐射树，$\mathcal P(j)$ 为
        根到母线 $j$ 的唯一路径，$\mathrm{sub}(j)$ 为 $j$ 的子树（含 $j$ 自身），
        $\ell_{ij}=|I_{ij}|^2$ 为支路电流幅值平方。
\end{itemize}

\subsection{方法谱系：不动点迭代视角下的经典方法}\label{sec:pf-th-cls-lineage}

潮流计算的数字化历史几乎与数字计算机同步。Ward 与 Hale 于 1956 年给出首个
实用的导纳阵迭代法~\cite{cls:wardhale}；随后出现了以阻抗阵 $Z_{bus}=Y_{bus}^{-1}$
为核心的一类方法（Z-bus 法），其系统整理见 Stagg 与 El-Abiad~\cite{cls:stagg}；
Tinney 与 Hart 于 1967 年确立 Newton 法加稀疏排序分解的主流地位~\cite{cls:tinney}；
Stott 与 Alsa\c{c} 于 1974 年提出快速解耦潮流~\cite{cls:stott}。Stott 的
综述~\cite{cls:stottreview} 对这一谱系有经典评述。本节用统一的\emph{不动点迭代}
语言重述前 Newton 时代方法，以看清其数学结构与被取代的原因。

\begin{definition}[Y-bus Gauss--Seidel 潮流迭代]\label{sec:pf-th-cls-def-gs}
把节点方程 $I=Y_{bus}V$ 与注入约束 $S_i^{\mathrm{spec}}=V_i\overline{I_i}$
联立，对每个 PQ 母线解出
\begin{equation}
V_i^{(k+1)}=\frac{1}{Y_{ii}}\Biggl(\frac{\overline{S_i^{\mathrm{spec}}}}{\overline{V_i^{(k)}}}
-\sum_{j<i}Y_{ij}V_j^{(k+1)}-\sum_{j>i}Y_{ij}V_j^{(k)}\Biggr),
\label{sec:pf-th-cls-eq:gs}
\end{equation}
即把潮流问题写成电压向量上的不动点方程 $V=T(V)$ 并做 Gauss--Seidel 型分量迭代。
PV 母线先由 $Q$ 方程解出无功再校正幅值。Z-bus 法把式~\eqref{sec:pf-th-cls-eq:gs}
右端的邻接求和换成 $Z_{bus}$ 与注入电流的乘积，收敛性更好但每步需作用稠密
$Z_{bus}$（或其隐式因子）。
\end{definition}

式~\eqref{sec:pf-th-cls-eq:gs} 中电流项 $\overline{S}/\overline{V}$ 使 $T$ 成为
\emph{非线性}映射；其收敛性分析的第一步是看线性化的骨架——对线性方程组
$YV=I$ 的 Gauss--Seidel 迭代，下述经典结论给出充分收敛条件。

\begin{theorem}[严格对角占优线性系统的 Gauss--Seidel 收敛性]\label{sec:pf-th-cls-thm-gs}
设 $A\in\mathbb C^{n\times n}$ 严格对角占优（每行
$|a_{ii}|>\sum_{j\ne i}|a_{ij}|$），则对任意初值，求解 $Ax=b$ 的
Gauss--Seidel 迭代收敛，迭代矩阵 $T_{GS}=-(D+L)^{-1}U$ 满足 $\rho(T_{GS})<1$
（其中 $A=D+L+U$ 为对角/严格下三角/严格上三角分解）。
\end{theorem}

\begin{proof}
设 $\lambda$ 为 $T_{GS}$ 的任一特征值，$x\ne0$ 为相应特征向量：
$-(D+L)^{-1}Ux=\lambda x$，即 $M(\lambda)x=0$，其中
$M(\lambda)=\lambda(D+L)+U$。$M(\lambda)$ 的元素为
$m_{ii}=\lambda a_{ii}$，$m_{ij}=\lambda a_{ij}$（$j<i$），$m_{ij}=a_{ij}$
（$j>i$）。设 $|\lambda|\ge1$，则对每行 $i$，
\begin{equation*}
|m_{ii}|=|\lambda|\,|a_{ii}|>|\lambda|\sum_{j\ne i}|a_{ij}|
\ge \sum_{j<i}|\lambda a_{ij}|+\sum_{j>i}|a_{ij}|=\sum_{j\ne i}|m_{ij}|,
\end{equation*}
即 $M(\lambda)$ 仍严格对角占优。严格对角占优矩阵非奇异（否则由
Gershgorin 圆盘定理存在含原点的圆盘，与对角占优矛盾），故 $M(\lambda)x=0$
只有零解，与 $x\ne0$ 矛盾。因此 $T_{GS}$ 的所有特征值满足 $|\lambda|<1$，
即 $\rho(T_{GS})<1$，迭代误差 $e^{(k)}=T_{GS}^{\,k}e^{(0)}\to0$。
\end{proof}

\begin{remark}[从线性骨架到非线性潮流]\label{sec:pf-th-cls-rem-gslim}
电力网络的 $Y_{bus}$ 通常弱对角占优而非严格占优（联络线、充电电容、三绕组
变压器等都会削弱占优性），且式~\eqref{sec:pf-th-cls-eq:gs} 的非线性项
$\overline{S}/\overline{V}$ 在电压偏低时放大局部 Lipschitz 常数，使收敛域收缩、
收敛速度随系统规模恶化。Gauss--Seidel 类方法只具有\emph{线性}（几何）收敛速度，
且渐近常数接近 1；这正是 Newton 法（局部二次收敛）~\cite{cls:tinney} 与其
近似版 FDPF~\cite{cls:stott} 取而代之的根本原因。作为不动点方法的价值在于结构
透明：它把"潮流 = 求 $V=T(V)$ 的不动点"这一观点留给了后来的前推回代
（第~\ref{sec:pf-th-cls-distflow}~节）与 HELM 等方法。
\end{remark}

\begin{gapnote}
Y-bus Gauss--Seidel 与 Z-bus 阻抗阵法在平台 \srcpath{src/power\_flow/} 中
\textbf{均无实现}（求解器谱系以 Newton 核、FDPF、直流、BFS 为主，
见实现层"架构总览"章）；本节谱系内容仅作理论定位。
\end{gapnote}

\subsection{快速解耦潮流：解耦原理推导与 XB/BX 变体}\label{sec:pf-th-cls-fdpf}

\subsubsection{弱耦合假设的数学化}

极坐标 Newton 校正方程按 $(\theta,V)$ 分块（推导见本手册 Newton 理论章）：
\begin{equation}
\begin{pmatrix}\Delta P\\ \Delta Q\end{pmatrix}
=
\begin{pmatrix} H & N\\ J & L\end{pmatrix}
\begin{pmatrix}\Delta\theta\\ \Delta V\end{pmatrix},
\label{sec:pf-th-cls-eq:nrblock}
\end{equation}
其中对 $i\ne j$（记 $B_{ij}=\operatorname{Im}Y_{ij}$），
\begin{align}
H_{ij}&=\frac{\partial P_i}{\partial\theta_j}
 =V_iV_j\bigl(G_{ij}\sin\theta_{ij}-B_{ij}\cos\theta_{ij}\bigr),&
N_{ij}&=V_j\frac{\partial P_i}{\partial V_j}
 =V_iV_j\bigl(G_{ij}\cos\theta_{ij}+B_{ij}\sin\theta_{ij}\bigr),\notag\\
J_{ij}&=\frac{\partial Q_i}{\partial\theta_j}
 =-V_iV_j\bigl(G_{ij}\cos\theta_{ij}+B_{ij}\sin\theta_{ij}\bigr),&
L_{ij}&=V_j\frac{\partial Q_i}{\partial V_j}
 =V_iV_j\bigl(G_{ij}\sin\theta_{ij}-B_{ij}\cos\theta_{ij}\bigr).
\label{sec:pf-th-cls-eq:blocks}
\end{align}

\begin{definition}[FDPF 三条物理解耦假设]\label{sec:pf-th-cls-def-assump}
\textbf{(A1) 高 $X/R$}：网络支路满足 $|G_{ij}|\ll|B_{ij}|$；
\textbf{(A2) 小角度差}：运行点处 $|\theta_{ij}|$ 小，
$\sin\theta_{ij}\approx0$、$\cos\theta_{ij}\approx1$；
\textbf{(A3) 无功通量小}：对角处 $|Q_i|\ll|B_{ii}|V_i^2$
（即流经的无功相对母线短路电纳口径为小量）。
\end{definition}

在 (A1)--(A2) 下，式~\eqref{sec:pf-th-cls-eq:blocks} 的非对角元化为
\begin{equation}
H_{ij}\approx L_{ij}\approx -V_iV_jB_{ij},\qquad
N_{ij}\approx -J_{ij}\approx V_iV_jG_{ij}\approx0,
\label{sec:pf-th-cls-eq:offdiag}
\end{equation}
即 $P$--$\theta$ 与 $Q$--$V$ 两块\textbf{解耦}，交叉块 $N,J$ 消失。
对角元 $H_{ii}=-Q_i-B_{ii}V_i^2$、$L_{ii}=Q_i-B_{ii}V_i^2$ 在 (A3) 下
均 $\approx -B_{ii}V_i^2$。于是两块校正方程可统一写成
\begin{equation}
\Delta P\approx -\,\mathrm{diag}(V)\,B'\,\mathrm{diag}(V)\,\Delta\theta,
\qquad
\Delta Q\approx -\,\mathrm{diag}(V)\,B''\,\mathrm{diag}(V)\,\Delta V,
\end{equation}
其中 $B',B''$ 是由 $-B$ 派生的 Laplacian 型常量阵（对角为正）。
两端除以 $V_i$ 并把右端残留的幅值因子取为 $1$（即再叠加一条
$V\approx1$ 的\textbf{(A4)}），得 Stott--Alsa\c{c} 标准形~\cite{cls:stott}
\begin{equation}
\boxed{\;B'\,\Delta\theta=\frac{\Delta P}{V},\qquad B''\,\Delta V=\frac{\Delta Q}{V}\;}
\label{sec:pf-th-cls-eq:fdpf}
\end{equation}
两阵\textbf{常量化}，可在迭代前一次因子分解，每次迭代只做失配计算与两次
稀疏回代；收敛判据一般取
$\max(\lVert\Delta P/V\rVert_\infty,\lVert\Delta Q/V\rVert_\infty)<\varepsilon$。

\subsubsection{XB 与 BX 变体：变压器分接头归属的差异}

两条假设链在"如何处理变压器变比 $t$、线路充电与母线并联"上留下自由度，
由此产生两个经典变体~\cite{cls:stott,cls:vanamerongen}：

\begin{definition}[XB 与 BX 方案]\label{sec:pf-th-cls-def-xbbx}
\textbf{XB 方案}（Stott--Alsa\c{c} 原始版，亦称 FDXB）：
$B'$ 由纯串联电纳 $1/x$ 装配——忽略电阻、充电、母线并联，\textbf{变比取 $1$}、
相移取 $0$；$B''$ \textbf{保留}变比、充电与母线并联的影响（仍忽略电阻与相移）。
\textbf{BX 方案}：分接头归属对调——$B'$ \textbf{保留}变比（含移相的注入等效），
$B''$ 中\textbf{变比取 $1$}；其余构造相同。
\end{definition}

\begin{remark}[变体选择的经验结论]\label{sec:pf-th-cls-rem-xbbx}
van Amerongen~\cite{cls:vanamerongen} 的"通用版"FDPF 与
Monticelli--Garcia--Saavedra~\cite{cls:monticelli} 对假设链的系统推导与测试
表明：XB 方案在常规输电网（$R/X$ 小、分接头接近 1）上收敛最快；当网络
$R/X$ 偏高或分接头远离标称时，BX 方案的收敛域与稳健性明显更好。物理原因：
分接头同时影响有功与无功通道，把它放进 $B'$（BX）使 $P$--$\theta$ 子迭代
吸收了主要的对角缩放效应；而 $R/X$ 大直接破坏假设 (A1)，两个方案都会退化，
此时应回退全 Newton 或改用对 $R/X$ 不敏感的方法（配网见
第~\ref{sec:pf-th-cls-distflow}~节）。
\end{remark}

\subsubsection{收敛性：近似 Newton 的线性收敛定理}

FDPF 不再是精确 Newton，而是\textbf{常量近似雅可比的拟 Newton 迭代}
$x^{k+1}=x^k-A^{-1}F(x^k)$，其中 $A=\mathrm{diag}(-V B' V,\,-V B'' V)\big|_{V=1}$。
其收敛性由下述 Ostrowski 型定理刻画（一般形式见~\cite{cls:ortega}）。

\begin{theorem}[常量近似雅可比迭代的局部线性收敛]\label{sec:pf-th-cls-thm-fdpf}
设 $F:\mathbb R^m\to\mathbb R^m$ 在解 $x^*$ 的邻域内二阶连续可微，
$F(x^*)=0$，$J(x^*)$ 非奇异；取常量非奇异阵 $A$。若
\begin{equation}
\sigma:=\rho\bigl(I-A^{-1}J(x^*)\bigr)<1,
\end{equation}
则存在 $x^*$ 的邻域，使迭代 $x^{k+1}=x^k-A^{-1}F(x^k)$ 对其中任意初值收敛到
$x^*$，且收敛是线性的：对任意 $\varepsilon>0$ 存在范数与常数 $C$ 使
$\lVert x^{k+1}-x^*\rVert\le(\sigma+\varepsilon)\lVert x^k-x^*\rVert
+C\lVert x^k-x^*\rVert^2$。
\end{theorem}

\begin{proof}
记 $e_k=x^k-x^*$，$M=I-A^{-1}J(x^*)$。在 $x^*$ 处 Taylor 展开：
$F(x^k)=J(x^*)e_k+r_k$，其中 $\lVert r_k\rVert\le C_2\lVert e_k\rVert^2$
（$F$ 二阶连续可微）。于是
\begin{equation}
e_{k+1}=e_k-A^{-1}\bigl(J(x^*)e_k+r_k\bigr)=Me_k-A^{-1}r_k.
\label{sec:pf-th-cls-eq:errrec}
\end{equation}
由谱半径的连续性，对任意 $\varepsilon>0$ 可取 $\mathbb R^m$ 上的范数使
$\lVert M\rVert\le\sigma+\varepsilon=:\eta<1$（取 $\varepsilon<(1-\sigma)/2$）。
式~\eqref{sec:pf-th-cls-eq:errrec} 给出
$\lVert e_{k+1}\rVert\le\eta\lVert e_k\rVert+C_3\lVert e_k\rVert^2
=\bigl(\eta+C_3\lVert e_k\rVert\bigr)\lVert e_k\rVert$。
取 $\delta>0$ 使 $\eta+C_3\delta<1$；则当 $\lVert e_0\rVert<\delta$ 时归纳得
$\lVert e_k\rVert\le(\eta+C_3\delta)^k\lVert e_0\rVert\to0$，且渐近收敛因子
不超过 $\eta=\sigma+\varepsilon$。
\end{proof}

\begin{remark}[该定理对 FDPF 的实际含义]\label{sec:pf-th-cls-rem-fdpfcond}
条件 $\rho(I-A^{-1}J(x^*))<1$ 正是解耦假设 (A1)--(A4) 的\textbf{定量形式}：
假设越接近成立，$A$ 越逼近 $J(x^*)$，$\sigma$ 越接近 $0$、几何收敛越快；
$R/X$ 升高或角度差拉大时 $\sigma$ 趋近甚至超过 1，表现为迭代缓慢或发散。
与 Newton 的局部二次收敛相比，FDPF 用"线性收敛 + 每步 $O(\mathrm{nnz})$
回代 + 一次分解"换取单步成本的数量级下降，在适定输电网上的实践表现是
$2\sim5$ 次迭代达到工程容差~\cite{cls:stott,cls:monticelli}。同理，FDPF
\textbf{没有} Newton 的回退链/限值切换等全局化机制，其收敛完全依赖初值落在
定理~\ref{sec:pf-th-cls-thm-fdpf} 的邻域内。
\end{remark}

\begin{gapnote}
平台 FDPF（\srcpath{src/power\_flow/fdpf\_solver.cpp}）实现的是 \textbf{XB 方案}
（$B'$ 变比取 1、忽略电阻/充电/并联，$B''$ 保留变比与充电/并联，构型由
\srcpath{make\_bp\_flags}/\srcpath{make\_bpp\_flags} 控制）；\textbf{BX 方案未实现}。
实现层手册称该求解器为 FDXB，与定义~\ref{sec:pf-th-cls-def-xbbx} 的命名一致。
\end{gapnote}

\subsection{直流潮流：三条假设的严格推导与误差分析}\label{sec:pf-th-cls-dc}

\subsubsection{从 AC 支路方程到 $P=B\theta$}

支路 $k$（$i\!\to\!j$，串联导纳 $y_k=g_k+jb_k$，忽略充电）from 端有功潮流的
精确 AC 表达式为
\begin{equation}
P_{ij}=V_i^2 g_k-V_iV_j\bigl(g_k\cos\theta_{ij}+b_k\sin\theta_{ij}\bigr).
\label{sec:pf-th-cls-eq:acbranch}
\end{equation}

\begin{definition}[直流潮流三条假设]\label{sec:pf-th-cls-def-dc}
\textbf{(D1) 忽略电阻}：$r_k=0$，故 $g_k=0$、$b_k=-1/x_k$；
\textbf{(D2) 单位电压}：$V_i\equiv1$（标幺）；
\textbf{(D3) 小角度}：$\sin\theta_{ij}\approx\theta_{ij}$。
\end{definition}

逐条施加于式~\eqref{sec:pf-th-cls-eq:acbranch}：(D1) 消去 $g_k$ 项；
(D2) 置 $V_i=V_j=1$；得 $P_{ij}=-b_k\sin\theta_{ij}=\sin\theta_{ij}/x_k$；
(D3) 给出线性支路方程
\begin{equation}
P_{ij}=\frac{\theta_i-\theta_j}{x_k}.
\label{sec:pf-th-cls-eq:dcbranch}
\end{equation}
对全部支路求和得节点平衡 $\sum_j P_{ij}=P_i$，写成矩阵形式即
\begin{equation}
\boxed{\;B\,\theta=P\;}
\label{sec:pf-th-cls-eq:dcpf}
\end{equation}
其中 $B$ 为以 $1/x_k$ 为权的\textbf{加权图 Laplacian}：$B_{ij}=-1/x_k$（$i\ne j$，
有支路相连时），$B_{ii}=\sum_{k\ni i}1/x_k$。剔除 slack 行/列并把 slack 相角
（通常取 $0$）的贡献移到右端，得降阶方程
$B_{\mathrm{red}}\theta_{\mathrm{red}}=P_{\mathrm{red}}-B_{\mathrm{red,slack}}\theta_{\mathrm{slack}}$。
移相变压器（相移 $\varphi_k$）在式~\eqref{sec:pf-th-cls-eq:dcbranch} 中表现为
$\theta_{ij}\to\theta_{ij}-\varphi_k$，装配时等价为 from/to 端一对反向注入
$\mp b_k\varphi_k$。支路潮流回算
$P_{f,k}=(\theta_i-\theta_j-\varphi_k)/x_k$。

\subsubsection{可解性：接地 Laplacian 的正定性}

\begin{theorem}[降阶 $B$ 阵的对称正定性]\label{sec:pf-th-cls-thm-dc}
设 AC 网络连通，全部在运支路 $x_k>0$，slack 集合非空。则剔除全部 slack
行/列后的降阶阵 $B_{\mathrm{red}}$ 对称正定，直流潮流方程
\eqref{sec:pf-th-cls-eq:dcpf} 有唯一解 $\theta_{\mathrm{red}}$。
\end{theorem}

\begin{proof}
对称性由构造显然。对任意 $\theta\in\mathbb R^{n}$，二次型可写成支路平方和：
\begin{equation}
\theta^{\top}B\theta=\sum_{k=(i,j)}\frac{1}{x_k}\bigl(\theta_i-\theta_j\bigr)^2\ge0.
\label{sec:pf-th-cls-eq:quad}
\end{equation}
若 $\theta^{\top}B\theta=0$，则每条支路两端相角相等；网络连通推出
$\theta$ 为常向量。今限制到非 slack 子空间：若
$\theta_{\mathrm{red}}^{\top}B_{\mathrm{red}}\theta_{\mathrm{red}}=0$，把
$\theta_{\mathrm{red}}$ 以 slack 分量取 $0$ 补全为 $\theta$，由
式~\eqref{sec:pf-th-cls-eq:quad} 仍有 $\theta^{\top}B\theta=0$（slack 相角取 0
时含 slack 的支路项为 $(\theta_i-0)^2$，不变号），故 $\theta$ 为常向量；
而 slack 分量为 $0$，所以常数为 $0$，即 $\theta_{\mathrm{red}}=0$。
因此 $B_{\mathrm{red}}$ 正定，必非奇异，解存在唯一。
\end{proof}

\begin{remark}\label{sec:pf-th-cls-rem-dcslack}
定理~\ref{sec:pf-th-cls-thm-dc} 说明：直流潮流的适定性只需要"连通 + 至少一个
相角参考 + 正电抗"，与运行点无关——这是它作为筛选/规划工具的鲁棒性来源。
代价是信息损失：$|V|$、$Q$、网损全部不输出，slack 母线自动吸收全系统不平衡
（在无损模型中恰为 $\sum P_i^{\mathrm{spec}}$ 的反号）。
\end{remark}

\subsubsection{线性化误差分析}

保留电阻与电压幅值后，直流模型相对式~\eqref{sec:pf-th-cls-eq:acbranch} 的误差
可以逐项界定。下述命题给出 $V\equiv1$ 情形下的干净分解。

\begin{proposition}[直流潮流支路线性化误差界]\label{sec:pf-th-cls-prop-dcerr}
设支路 $k$ 阻抗 $z_k=r_k+jx_k$（$x_k>0$），$V_i=V_j=1$，记
$\theta=\theta_{ij}$，$\rho_k=r_k/x_k$。则精确有功
$P_{ij}=g_k(1-\cos\theta)+\frac{x_k}{r_k^2+x_k^2}\sin\theta$，
与直流近似 $\theta/x_k$ 的误差 $\varepsilon_k=P_{ij}-\theta/x_k$ 满足
\begin{equation}
|\varepsilon_k|
\le \frac{g_k\theta^2}{2}
+\frac{\rho_k^2}{x_k(1+\rho_k^2)}\,|\theta|
+\frac{|\theta|^3}{6x_k}.
\label{sec:pf-th-cls-eq:dcerr}
\end{equation}
三项分别对应\textbf{损耗项}（$g(1-\cos\theta)$，二阶小量但恒正、系统性偏高）、
\textbf{有效电抗项}（$1/x_{\mathrm{eff}}=x_k/(r_k^2+x_k^2)$ 与 $1/x_k$ 之差，
$\rho_k^2$ 量级）与\textbf{正弦线性化项}（三阶小量）。
\end{proposition}

\begin{proof}
由 $y_k=1/z_k=(r_k-jx_k)/(r_k^2+x_k^2)$，得 $g_k=r_k/(r_k^2+x_k^2)$、
$b_k=-x_k/(r_k^2+x_k^2)$，代入式~\eqref{sec:pf-th-cls-eq:acbranch}
（$V_i=V_j=1$）即得精确表达式。于是
\begin{equation*}
\varepsilon_k=g_k(1-\cos\theta)
+\sin\theta\Bigl(\frac{x_k}{r_k^2+x_k^2}-\frac{1}{x_k}\Bigr)
+\frac{\sin\theta-\theta}{x_k}.
\end{equation*}
逐项取绝对值：$0\le1-\cos\theta\le\theta^2/2$；
$\bigl|\frac{x_k}{r_k^2+x_k^2}-\frac{1}{x_k}\bigr|
=\frac{r_k^2}{x_k(r_k^2+x_k^2)}=\frac{\rho_k^2}{x_k(1+\rho_k^2)}$；
$|\sin\theta-\theta|\le|\theta|^3/6$（Taylor 余项的交错级数估计）。
三角不等式合并即得式~\eqref{sec:pf-th-cls-eq:dcerr}。
\end{proof}

\begin{remark}[误差的量级解读]\label{sec:pf-th-cls-rem-dcerr}
输电网典型 $\rho_k\le0.3$、$|\theta|\le0.5$~rad：有效电抗项
$\sim\rho_k^2|\theta|/x_k\le4.5\%$（相对 $|\theta|/x_k$），损耗项与正弦项更小；
这与工程文献中"直流潮流支路有功误差约百分之几"的经验一致~\cite{cls:stottreview}。
但对\textbf{高 $R/X$ 配网}（$\rho_k$ 可超 1），有效电抗项达 $O(\rho_k^2)$，
式~\eqref{sec:pf-th-cls-eq:dcerr} 失去实用意义——直流潮流不适用于配网，
与 FDPF 的失效机理同源（假设 (A1)/(D1)）。
\end{remark}

\subsubsection{损耗近似}

直流模型自身无损（$g=0$ 隐含 $\sum P_i^{\mathrm{spec}}=0$ 才可精确平衡），
工程中需要网损估计时标准的做法是\textbf{事后二次近似}：由精确支路损耗
$P_{\mathrm{loss},k}=g_k|V_i-V_j|^2=g_k\bigl(V_i^2+V_j^2-2V_iV_j\cos\theta_{ij}\bigr)$，
施加 (D2) 并展开 $\cos$ 到二阶，得
\begin{equation}
P_{\mathrm{loss},k}\approx g_k\,\theta_{ij}^2=\frac{r_k}{r_k^2+x_k^2}\,\theta_{ij}^2
\approx r_k\Bigl(\frac{\theta_{ij}}{x_k}\Bigr)^2
=r_k\,\frac{P_{f,k}^2}{1},
\label{sec:pf-th-cls-eq:loss}
\end{equation}
即"损耗 $\approx r\times$（直流潮流）$^2$"的二次口径。进一步可把
$\sum_k P_{\mathrm{loss},k}/2$ 按比例回加到负荷并重解（损耗修正直流潮流），
使 slack 吸收量逼近真实网损。

\begin{gapnote}
平台直流求解器 \srcpath{solve\_ac\_linearized\_dc} 输出仅为相角与无损支路潮流；
式~\eqref{sec:pf-th-cls-eq:loss} 的\textbf{损耗事后估计与损耗修正重解均未实现}
（结果结构 \texttt{ACLinearizedDCResult} 无损耗字段）。
\end{gapnote}

\subsection{直流灵敏度：转移分布因子与线路开断分布因子}\label{sec:pf-th-cls-ptdf}

直流模型的线性性使支路潮流对节点注入、以及对线路开断的\emph{灵敏度}具有闭式矩阵
表达。这两类因子——PTDF 与 LODF——是安全分析、市场出清与 N-1 校核的基础工具，
其价值在于\textbf{只依赖网络拓扑与电抗、与运行点无关}，可离线预计算。

\subsubsection{PTDF：注入转移分布因子}
设剔除 slack 后的降阶节点电纳阵 $B\in\mathbb R^{(n-1)\times(n-1)}$
（式~\eqref{sec:pf-th-cls-eq:dcpf}，$B=A^{\top}B_dA$），支路-节点关联阵
$A\in\mathbb R^{b\times(n-1)}$（第 $l$ 行对应支路 $l=(i\!\to\!j)$，$i$ 列为 $+1$、
$j$ 列为 $-1$，slack 列剔除），支路电纳对角阵 $B_d=\mathrm{diag}(1/x_l)$。支路有功
$f_l=(\theta_i-\theta_j)/x_l$ 即 $f=B_dA\theta$，结合 $\theta=B^{-1}P$ 得

\begin{definition}[PTDF/ISF 矩阵]\label{def:pf-th-cls-ptdf}
\begin{equation}
f=\Phi\,P,\qquad \Phi:=B_dA\,B^{-1}\in\mathbb R^{b\times(n-1)} .
\label{sec:pf-th-cls-eq:ptdf}
\end{equation}
$\Phi_{lk}$ 是“节点 $k$ 单位注入、slack 回收”引起的支路 $l$ 潮流变化，称\emph{注入
转移分布因子}（injection shift factor, ISF），亦称\emph{功率转移分布因子}（PTDF）。
\end{definition}

\begin{proposition}[点对点交易的支路灵敏度]\label{prop:pf-th-cls-ptdf-pair}
从节点 $s$ 注入、节点 $t$ 抽取的单位交易（$P_s{=}{+}1,\,P_t{=}{-}1$）引起支路 $l$ 的
潮流变化为 $\Phi_{ls}-\Phi_{lt}$；特别地，沿支路 $m=(p\!\to\!q)$ 两端的单位对流引起
支路 $l$ 的变化记 $\rho_{lm}:=\Phi_{lp}-\Phi_{lq}$（slack 端分量取零）。
\end{proposition}
（由式~\eqref{sec:pf-th-cls-eq:ptdf} 对注入的线性叠加即得。）

\subsubsection{LODF：线路开断分布因子}
\begin{theorem}[Wood--Wollenberg LODF 公式]\label{thm:pf-th-cls-lodf}
设支路 $m$ 开断前潮流为 $f_m^0$，且开断不致系统解列（$\rho_{mm}\ne1$）。则开断后
支路 $l$（$l\ne m$）的潮流增量为
\begin{equation}
\Delta f_l=\psi_{lm}\,f_m^0,\qquad
\psi_{lm}=\frac{\rho_{lm}}{1-\rho_{mm}},
\label{sec:pf-th-cls-eq:lodf}
\end{equation}
$\rho_{lm}=\Phi_{lp}-\Phi_{lq}$、$\rho_{mm}=\Phi_{mp}-\Phi_{mq}$（$m=(p\!\to\!q)$），
$\psi_{lm}$ 称\emph{线路开断分布因子}（LODF）。
\end{theorem}
\begin{proof}
用补偿定理：在\emph{完整}网络上叠加一对沿 $m$ 端点的注入（$p$ 注入 $\Delta$、$q$
抽取 $\Delta$），以其等效替代被开断支路的电流。该注入使支路 $m$ 上产生附加流
$\rho_{mm}\Delta$，使任意支路 $l$ 产生 $\rho_{lm}\Delta$。开断的自洽条件是：这对注入
所“代表”的幻象电流 $\Delta$ 恰等于叠加后支路 $m$ 的总流，即
$\Delta=f_m^0+\rho_{mm}\Delta$，解得 $\Delta=f_m^0/(1-\rho_{mm})$。于是任意
$l\ne m$ 的增量 $\Delta f_l=\rho_{lm}\Delta=\dfrac{\rho_{lm}}{1-\rho_{mm}}f_m^0$。
$\rho_{mm}\to1$ 对应开断使 $m$ 成为割边、网络解列，LODF 发散。
\end{proof}

\begin{remark}[拓扑不变性与用途]\label{rem:pf-th-cls-lodf}
$\Phi$ 与 $\psi$ 只由 $A$、$B_d$（拓扑与电抗）决定，与负荷、发电水平无关，故可离线
预计算、在线以一次矩阵-向量积完成全部注入或单一开断的支路潮流筛查——这是直流
线性化在安全约束调度、N-1 校核与市场阻塞定价中被广泛采用的根本原因；其代价是
继承直流模型的全部近似（命题~\ref{sec:pf-th-cls-prop-dcerr}）。多重开断需对
$\psi$ 作分块推广（同时开断的支路间存在耦合）。
\end{remark}

\subsection{辐射状配电网潮流：DistFlow 模型与前推回代}\label{sec:pf-th-cls-distflow}

\subsubsection{DistFlow 支路潮流模型}

配电网的结构性差异——辐射拓扑、高 $R/X$、显著电压降落——使输电网方法的前提
（(A1)、(D1)--(D3)）不再成立，催生了以支路变量为核心的模型体系。
Baran 与 Wu~\cite{cls:baranwu} 的 DistFlow 模型是其中的标准形式。

\begin{definition}[DistFlow 支路潮流方程]\label{sec:pf-th-cls-def-distflow}
设辐射树 $\mathcal T$ 的支路一律按远离根的方向定向，支路 $(i,j)$ 阻抗
$z_{ij}=r_{ij}+jx_{ij}$，记 $P_{ij},Q_{ij}$ 为从 $i$ 流向 $j$ 的有功/无功，
$\ell_{ij}=|I_{ij}|^2$，$v_i=|V_i|^2$，$p_j,q_j$ 为母线 $j$ 的净负荷
（含本地并联的 $|V|^2$ 项）。由 $V_j=V_i-z_{ij}I_{ij}$ 与
$S_{ij}=V_i\overline{I_{ij}}$ 两端取模平方并分离实虚部，得递推方程组
\begin{align}
P_{ij}-r_{ij}\ell_{ij}-p_j&=\sum_{k\in\mathcal C(j)}P_{jk},
\label{sec:pf-th-cls-eq:dfp}\\
Q_{ij}-x_{ij}\ell_{ij}-q_j&=\sum_{k\in\mathcal C(j)}Q_{jk},
\label{sec:pf-th-cls-eq:dfq}\\
v_j&=v_i-2\bigl(r_{ij}P_{ij}+x_{ij}Q_{ij}\bigr)+\bigl(r_{ij}^2+x_{ij}^2\bigr)\ell_{ij},
\label{sec:pf-th-cls-eq:dfv}\\
\ell_{ij}\,v_i&=P_{ij}^2+Q_{ij}^2,
\label{sec:pf-th-cls-eq:dfl}
\end{align}
$\mathcal C(j)$ 为 $j$ 的子节点集。方程
\eqref{sec:pf-th-cls-eq:dfp}--\eqref{sec:pf-th-cls-eq:dfq} 是含损功率平衡，
\eqref{sec:pf-th-cls-eq:dfv} 是电压降落恒等式，\eqref{sec:pf-th-cls-eq:dfl}
是支路潮流与电流的耦合约束（唯一的非线性，且为凸二阶锥形式）。
略去损耗项 $r\ell$、$x\ell$ 与 $(r^2+x^2)\ell$ 得\textbf{ LinDistFlow}
（线性 DistFlow）模型，广泛用于配网优化建模~\cite{cls:baranwu}。
\end{definition}

\begin{remark}[存在唯一性]\label{sec:pf-th-cls-rem-chiang}
Chiang 与 Baran~\cite{cls:chiangbaran} 证明：对辐射状配网，在"负荷不过分重"
的可操作条件下，潮流方程的可行解在物理有意义的电压区域内\textbf{存在且唯一}，
且可用不动点框架构造。这一结果把配网潮流与输电网多解问题区分开：辐射 +
单电源 + 适度负荷保证了良态性，是前推回代可靠工作的理论底座。
\end{remark}

\subsubsection{前推回代：作为不动点迭代的推导}

前推回代（Backward--Forward Sweep, BFS；Shirmohammadi
等~\cite{cls:shirmohammadi}）不显式组装 DistFlow 方程组，而是把辐射树结构
直接变成两遍扫描的不动点迭代。

\begin{definition}[BFS 扫描映射]\label{sec:pf-th-cls-def-bfs}
给定电压剖面 $V=(V_1,\dots,V_n)$（根 $V_0$ 固定为 slack 电压），
\textbf{回推}（叶$\to$根）按
$I_{ij}=\sum_{m\in\mathrm{sub}(j)}\overline{S_m^{\mathrm{net}}/V_m}$
汇总各支路电流；\textbf{前推}（根$\to$叶）沿树按
$V'_j=V'_i-z_{ij}I_{ij}$ 更新电压。记一遍完整扫描为映射 $V'=T(V)$，
BFS 即 Picard 迭代 $V^{(k+1)}=T(V^{(k)})$，收敛判据
$\max_j|V_j^{(k+1)}-V_j^{(k)}|<\varepsilon$。
\end{definition}

展开前推递推：$V'_j=V_0-\sum_{l\in\mathcal P(j)}z_l I_l$。由于 $I_l$ 只依赖
$\mathrm{sub}$ 内母线电压，$T$ 是良定义的显式映射；其不动点精确满足
DistFlow 方程（回推保证
\eqref{sec:pf-th-cls-eq:dfp}--\eqref{sec:pf-th-cls-eq:dfq}，
前推保证 \eqref{sec:pf-th-cls-eq:dfv}，电流定义保证
\eqref{sec:pf-th-cls-eq:dfl}）。

\begin{theorem}[BFS 的压缩收敛性（充分条件）]\label{sec:pf-th-cls-thm-bfs}
设辐射网以母线 $0$ 为 slack，$V_0=1\angle0$ 固定；负荷恒功率，
$|S_j^{\mathrm{net}}|\le\sigma_j$，记 $\sigma_\Sigma=\sum_j\sigma_j$；支路阻抗
满足 $Z_{\max}:=\max_j\sum_{l\in\mathcal P(j)}|z_l|<\infty$（根到任一母线的
路径阻抗上界）。取 $v_{\min}\in(0,1)$ 并限制电压剖面于
$\mathcal D=\{V:\ |V_j|\ge v_{\min}\ \forall j\}$。若
\begin{equation}
L:=\frac{Z_{\max}\,\sigma_\Sigma}{v_{\min}^2}<1
\qquad\text{且}\qquad
1-\frac{Z_{\max}\sigma_\Sigma}{v_{\min}}\ge v_{\min},
\label{sec:pf-th-cls-eq:bfsc}
\end{equation}
则 $T$ 在 $\mathcal D$ 上是压缩映射（$\lVert T(V)-T(W)\rVert_\infty\le
L\lVert V-W\rVert_\infty$）且 $T(\mathcal D)\subseteq\mathcal D$；
由 Banach 不动点定理，$T$ 在 $\mathcal D$ 内有\textbf{唯一}不动点，
BFS 从 $\mathcal D$ 内任意初值出发以速率 $L$ 几何收敛。
\end{theorem}

\begin{proof}
\textbf{第一步（回推的 Lipschitz 界）。}对
$I_l=\sum_{m\in\mathrm{sub}(l)}\overline{S_m/V_m}$，在 $\mathcal D$ 上
\begin{equation*}
\Bigl|\frac{\overline{S_m}}{\overline{V_m}}-\frac{\overline{S_m}}{\overline{W_m}}\Bigr|
=\frac{|S_m|\,|V_m-W_m|}{|V_m|\,|W_m|}
\le\frac{\sigma_m}{v_{\min}^2}\,\lVert V-W\rVert_\infty,
\end{equation*}
对子树求和得
$|I_l(V)-I_l(W)|\le\bigl(\sigma_\Sigma/v_{\min}^2\bigr)\lVert V-W\rVert_\infty$。
\textbf{第二步（前推传播）。}由 $T(V)_j=V_0-\sum_{l\in\mathcal P(j)}z_l I_l(V)$，
\begin{equation*}
|T(V)_j-T(W)_j|\le\sum_{l\in\mathcal P(j)}|z_l|\,
\frac{\sigma_\Sigma}{v_{\min}^2}\lVert V-W\rVert_\infty
\le L\,\lVert V-W\rVert_\infty,
\end{equation*}
对 $j$ 取上确界即压缩性。\textbf{第三步（自映射）。}对 $V\in\mathcal D$，
$|I_l(V)|\le\sigma_\Sigma/v_{\min}$，故
$|T(V)_j|\ge1-Z_{\max}\sigma_\Sigma/v_{\min}\ge v_{\min}$
（式~\eqref{sec:pf-th-cls-eq:bfsc} 第二式），即 $T(V)\in\mathcal D$。
$\mathcal D$ 是有限维空间中的闭集（完备度量空间），Banach 定理给出
唯一不动点与几何收敛 $\lVert V^{(k)}-V^*\rVert_\infty\le
\frac{L^k}{1-L}\lVert V^{(1)}-V^{(0)}\rVert_\infty$。
\end{proof}

\begin{remark}[条件 \eqref{sec:pf-th-cls-eq:bfsc} 的物理读法]\label{sec:pf-th-cls-rem-bfsc}
$Z_{\max}\sigma_\Sigma$ 是"路径阻抗 $\times$ 总负荷"——即\textbf{短路容量相对
负荷规模的裕度}：条件要求馈线足够"强"（短路容量大、负荷轻、电压不贴近崩塌点
$v_{\min}$ 过小）。这与工程直觉一致：重载长馈线末端电压低，BFS 迭代变慢甚至
失效；也解释了为何实现必须用拓扑 fail-close 与残差监控兜底，而不能依赖先验
条件检验（$L<1$ 只是充分条件，实践中 $L$ 略超 1 仍常收敛）。
恒阻抗/恒电流负荷使 $T$ 更接近线性映射（注入电流对 $V$ 的敏感度降低），
压缩条件更宽松——这就是 ZIP 负荷有助于 BFS 收敛的理论原因。
\end{remark}

\begin{remark}[高 $R/X$ 弱网特性与方法选择]\label{sec:pf-th-cls-rem-rx}
配网 $R/X$ 常大于 1：(i) FDPF/直流的假设 (A1)/(D1) 被破坏
（命题~\ref{sec:pf-th-cls-prop-dcerr} 的有效电抗项达 $O(\rho_k^2)$）；
(ii) $P$--$V$、$Q$--$\theta$ 交叉耦合增强，解耦迭代发散；
(iii) 电压降落项 $2(rP+xQ)$ 中电阻贡献占主导，$|V|$ 对有功潮流敏感。
BFS 与 DistFlow 不显式使用任何弱耦合假设（式
\eqref{sec:pf-th-cls-eq:dfp}--\eqref{sec:pf-th-cls-eq:dfl} 是精确恒等式），
对 $R/X$ 天然免疫，配合辐射拓扑的 $O(n)$ 扫描复杂度，成为配网正序潮流的
标准选择；环网/弱环网则需补偿法~\cite{cls:shirmohammadi} 或三相 Newton。
\end{remark}

\begin{gapnote}
平台实现的是\textbf{精确 AC 支路模型的 BFS}
（\srcpath{run\_bfs\_sweep}，含变比折算与 ZIP 电压依赖），收敛点与
DistFlow 方程组等价，但\textbf{没有}独立的 DistFlow 方程组求解器；
LinDistFlow 仅以 MIP 建模层形式存在于 power\_models 模块
（\srcpath{src/power\_models/lindistflow\_builder.cpp:solve\_lindistflow}），
不属于潮流求解路径。定理~\ref{sec:pf-th-cls-thm-bfs} 的压缩条件
\eqref{sec:pf-th-cls-eq:bfsc} 亦\textbf{未}作为先验检验实现——实现以辐射/连通
fail-close 检查加 $\max|\Delta V|$ 收敛判据代替。
\end{gapnote}

\subsection{分布式松弛与自适应孤岛的理论定位}\label{sec:pf-th-cls-slack}

\subsubsection{松弛母线的自由度账与物理意义}

$n$ 母线 AC 潮流每节点有四个量 $P,Q,|V|,\theta$、两条规范方程（注入定义），
因此每母线必须指定两个量；但全系统的\textbf{网损在求解前未知}，若所有母线的
$P$ 都事先钉死，功率平衡 $\sum P_i=\text{损耗}$ 一般无解。松弛（slack）母线的
本质是\textbf{开放其 $P,Q$ 两个自由度}以吸收先验未知的不平衡（主要是网损），
同时提供相角参考——$\theta$ 只在连通分量内有相对意义，必须钉住一个零点
（自由度账：$n$ 个 PQ/PV 母线贡献 $2n-2$ 个失配方程，对应 $\theta$ 与 $|V|$
的 $2n-2$ 个未知量，slack 的 $\theta,|V|$ 钉值）。

\begin{definition}[分布式松弛与参与因子]\label{sec:pf-th-cls-def-dslack}
把单一 slack 吸收的不平衡 $\Delta P_\Sigma$ 按归一化\textbf{参与因子}
$\alpha_i\ge0$、$\sum_i\alpha_i=1$ 分摊到一组可调电源母线
$\mathcal G$：$P_{g,i}\leftarrow P_{g,i}+\alpha_i\Delta P_\Sigma$。
因子的常见生成口径：按容量（$\alpha_i\propto P_{g,i}^{\max}$）、
按下垂系数（$\alpha_i\propto1/R_i$，对应一次调频的稳态份额）或均分。
\end{definition}

\begin{remark}[slack 分摊的物理意义与两种数学口径]\label{sec:pf-th-cls-rem-dslack}
单 slack 模型隐含"一台无穷大机组瞬时平衡全网"，在网损分摊、AGC 与经济调度的
衔接上不真实；分布式松弛把不平衡按参与因子映射到多台机组，模拟调速器
稳态分配。\emph{数学上}有两种口径：\textbf{方程内口径}把
$P_{g,i}=P_{g,i}^0+\alpha_i\,\pi$（$\pi$ 为增广未知的不平衡量）直接嵌入
Newton 失配方程同步求解，结果自洽（分摊后网损变化已被迭代吸收）；
\textbf{事后口径}先用单 slack 解出基准潮流，再把 slack 吸收量按 $\alpha$ 外部分摊。
事后口径忽略分摊引起的网损二阶变化（$\Delta P_\Sigma$ 改变潮流分布进而改变
损耗），属于一阶近似；要自洽需把分摊值写回后\textbf{重解}直至残余失配收敛。
限值处理（机组 $P_{\min}/P_{\max}$ 越限钉限再分摊剩余）使口径进一步成为
带投影的迭代。
\end{remark}

\subsubsection{方程内精确分布式松弛的加边 Newton 系统}\label{sec:pf-th-cls-augslack}
注记~\ref{sec:pf-th-cls-rem-dslack} 的“事后口径”只一阶近似分摊引起的网损变化。
\emph{方程内}口径把分摊量作为增广未知量与潮流方程联立，自洽吸收网损的二阶效应。

\begin{proposition}[加边分布式松弛系统]\label{prop:pf-th-cls-augslack}
引入标量未知 $\pi$（待分摊的总不平衡）；参与母线 $i\in\mathcal G$ 的有功注入取
$P_i^{\mathrm{spec}}(\pi)=P_i^0+\alpha_i\pi$（$\sum_{i\in\mathcal G}\alpha_i=1$）。
保留一个\emph{仅}提供角度与电压基准的参考母线（$\theta_{\mathrm{ref}}=0$、
$V_{\mathrm{ref}}$ 固定、其无功自由），并\textbf{恢复}参考母线的有功平衡方程纳入方程组。
则未知量 $(\theta_{\ne\mathrm{ref}},V_{m,\mathrm{PQ}},\pi)$ 与方程（全部母线有功平衡
$+$ PQ 母线无功平衡）个数相等，其 Newton 系统为\emph{加边}结构
\begin{equation}
\begin{bmatrix} J_{\mathrm{pf}} & \boldsymbol\alpha\\[2pt]
\boldsymbol c^{\top} & 0\end{bmatrix}
\begin{bmatrix}\Delta x\\ \Delta\pi\end{bmatrix}
=-\begin{bmatrix}F_{\mathrm{pf}}\\ F_{\mathrm{ref}}\end{bmatrix},
\label{sec:pf-th-cls-eq:augslack}
\end{equation}
其中 $J_{\mathrm{pf}}$ 为常规潮流雅可比，$\boldsymbol\alpha$ 为
$\partial F_{\mathrm{pf}}/\partial\pi$ 的参与因子列，$\boldsymbol c^{\top}$ 为新增
参考母线有功平衡对 $x$ 的偏导行。
\end{proposition}

\begin{remark}[自洽性与加边求解]\label{rem:pf-th-cls-augslack}
式~\eqref{sec:pf-th-cls-eq:augslack} 的解 $\pi^\star$ 是收敛点处各参与母线共同吸收的
总不平衡（对应全网损耗份额）；网损随分摊变化的二阶效应已在迭代内消化，正是事后
口径所缺。加边系统只比基潮流多一行一列，用块消元（Schur 补/矩阵求逆引理）可在\emph{一次}
基潮流分解上求解——$\Delta x=-J_{\mathrm{pf}}^{-1}(F_{\mathrm{pf}}+\boldsymbol\alpha\Delta\pi)$
代入底行得关于标量 $\Delta\pi$ 的方程，$\Delta\pi=-(F_{\mathrm{ref}}-\boldsymbol c^{\top}
J_{\mathrm{pf}}^{-1}F_{\mathrm{pf}})/(\boldsymbol c^{\top}J_{\mathrm{pf}}^{-1}\boldsymbol\alpha)$，
仅需两次以 $J_{\mathrm{pf}}$ 的回代，无需如“重复完整重解”那样多次重构与再分解。
多参考（多岛）时 $\pi$ 逐岛设立、$\alpha$ 逐岛归一。
\end{remark}

\begin{gapnote}
式~\eqref{sec:pf-th-cls-eq:augslack} 的方程内加边分布式松弛在平台中\textbf{未实现}
（实现为事后分摊与重复重解，见本节末“与实现的对应”表及其函数命名说明）。
\end{gapnote}

\subsubsection{自适应孤岛：连通分量、参考角与死岛}

拓扑分裂（故障开断、计划孤岛）使一个求解对象变成多个电气解耦的子问题。
其代数判据是图 Laplacian 的秩亏结构。

\begin{lemma}[Laplacian 零空间的维数 = 连通分量数]\label{sec:pf-th-cls-lem-lap}
设无向图有 $n$ 个顶点、$c$ 个连通分量，边权 $w_k>0$，$B$ 为其加权
Laplacian（$B_{ij}=-w_k$、$B_{ii}=\sum_{k\ni i}w_k$）。则
$\dim\ker B=c$，$\ker B$ 由各分量的示性向量张成。
\end{lemma}

\begin{proof}
由式~\eqref{sec:pf-th-cls-eq:quad} 的二次型展开
$x^{\top}Bx=\sum_{k=(i,j)}w_k(x_i-x_j)^2\ge0$：$Bx=0\Rightarrow x^{\top}Bx=0$
$\Rightarrow$ 每条边两端分量相等 $\Rightarrow$ $x$ 在每个连通分量上取常值。
反之每个分量上的常值向量显然在 $\ker B$ 中。故 $\ker B$ 恰为
$\{\text{各分量示性向量的线性组合}\}$，维数为 $c$。
\end{proof}

\begin{remark}[引理~\ref{sec:pf-th-cls-lem-lap} 对潮流求解的含义]\label{sec:pf-th-cls-rem-island}
$c$ 个孤岛的网络，其潮流雅可比（以及直流 $B$ 阵）天然有 $c$ 重零特征值：
每个孤岛都需要\textbf{各自}的相角参考与功率平衡机制，混用单一 slack 的
整体方程组必然奇异。\emph{自适应孤岛}的理论定位因此是：先分岛（图的连通分量
分解），把原问题分解为 $c$ 个独立的适定子问题，活岛各自指定/自选 slack 求解；
\textbf{无源岛}（无任何注入设备）在潮流框架内\textbf{无可行解}——其负荷无法满足、
电压无支撑，物理上对应失压停运，诚实的计算处理是作为\emph{死岛}以零电压
哨兵标记并排除出求解，而非强行给出无意义数值。这与定理
\ref{sec:pf-th-cls-thm-dc} 的适定性条件（连通 + 参考 + 正电抗）互为表里。
\end{remark}

\begin{gapnote}
方程内口径的精确分布式松弛（$\pi$ 作为增广未知量进入 Newton 失配）在平台中
\textbf{未实现}：\srcpath{DistributedSlackSolver::solve\_simplified} 是事后分摊
（一次基准 Newton + \srcpath{compute\_total\_participating\_mismatch\_pu} 复算
分摊），\srcpath{solve\_full\_jacobian} 是"重调度 + 反复完整重解 + 钉限重分摊"
的迭代逼近（兼容名 \texttt{full\_jacobian} 不表示增广 Jacobian）。
\end{gapnote}

\subsection{数值算例与基准}\label{sec:pf-th-cls-numexp}

本节数字由 \srcpath{tools/pf\_doc\_benchmark.cpp}（模式 \texttt{fdpf}/\texttt{dc}/\texttt{bfs}）
在 MATPOWER 算例上经平台公共求解 API 实测生成（平启动，Apple~M4，单线程），
迭代次数与误差为机器无关量。

\subsubsection{FDPF 与 Newton 的迭代数对比}
表~\ref{tab:pf-th-cls-fdpf} 给出统一 Newton 与 FDPF（XB 方案）达到 $\mathrm{tol}=10^{-8}$
所需迭代数。FDPF 迭代数（$11\sim44$）显著多于 Newton（$3\sim8$），且随系统解耦质量
（$R/X$、角度差）波动——这正是定理~\ref{sec:pf-th-cls-thm-fdpf} 的\emph{线性}收敛与
Newton \emph{二次}收敛之别；FDPF 以更多次但每次仅两回代（常量 $B'/B''$ 一次分解）
换取单步成本的数量级下降。
\begin{table}[H]\centering\footnotesize
\caption{FDPF 与统一 Newton 的迭代数（平启动，$\mathrm{tol}=10^{-8}$）}
\label{tab:pf-th-cls-fdpf}
\begin{tabular}{@{}lrrr@{}}
\toprule
算例 & 母线数 $n$ & Newton 迭代 & FDPF(XB) 迭代\\
\midrule
case9   & 9   & 5 & 11\\
case14  & 14  & 3 & 26\\
case30  & 30  & 4 & 27\\
case57  & 57  & 4 & 39\\
case118 & 118 & 7 & 15\\
case300 & 300 & 8 & 44\\
\bottomrule
\end{tabular}
\end{table}

\subsubsection{直流潮流的支路有功误差}
表~\ref{tab:pf-th-cls-dc} 给出直流线性化支路有功与全交流解的偏差（逐支路 from 端有功）。
平均误差为 $0.4\sim11$~MW（相对各算例总负荷约 $0.05\%\sim0.5\%$），但重载支路的最大误差
可达数十至数百 MW，随规模/负荷压力上升——与命题~\ref{sec:pf-th-cls-prop-dcerr} 的误差
界（损耗项 $\propto\theta^2$、正弦线性化项 $\propto\theta^3$）一致：角度差越大，线性化
误差越显著。直流适合把握整体潮流格局，不宜用于重载支路的精确定量。
\begin{table}[H]\centering\footnotesize
\caption{直流线性化相对全交流的支路有功误差（MW）}
\label{tab:pf-th-cls-dc}
\begin{tabular}{@{}lrrrrr@{}}
\toprule
算例 & 母线数 & 支路数 & 最大误差 & 平均误差 & 均方根误差\\
\midrule
case9   & 9   & 9   & $4.6$   & $1.25$ & $1.95$\\
case14  & 14  & 20  & $9.0$   & $1.25$ & $2.42$\\
case30  & 30  & 41  & $1.7$   & $0.39$ & $0.53$\\
case57  & 57  & 80  & $9.6$   & $1.57$ & $2.45$\\
case118 & 118 & 186 & $59.4$  & $3.60$ & $7.46$\\
case300 & 300 & 411 & $409.5$ & $10.6$ & $29.8$\\
\bottomrule
\end{tabular}
\end{table}

\subsubsection{前推回代在配电馈线上的表现}
表~\ref{tab:pf-th-cls-bfs} 在四个辐射配电算例上对比 BFS 与统一 Newton。两法最低母线电压
一致到 $10^{-9}$（交叉校核），BFS 迭代数（$7\sim9$）略多于 Newton（$4\sim5$），但每次
仅 $O(n)$ 次树上扫描、无雅可比装配与分解；馈线末端电压低至 $0.87\sim0.93$~pu，体现
高 $R/X$ 长馈线的显著压降。BFS 对 $R/X$ 天然免疫（式~\eqref{sec:pf-th-cls-eq:dfp}--%
\eqref{sec:pf-th-cls-eq:dfl} 为精确恒等式），是注记~\ref{sec:pf-th-cls-rem-rx} 所述配网
首选方法的实测印证。
\begin{table}[H]\centering\footnotesize
\caption{前推回代（BFS）与统一 Newton 在辐射配电算例上的对比（$\mathrm{tol}=10^{-8}$）}
\label{tab:pf-th-cls-bfs}
\begin{tabular}{@{}lrrrrrr@{}}
\toprule
算例 & 母线数 & BFS 迭代 & BFS $\min|V|$ & Newton 迭代 & Newton $\min|V|$ & $\min|V|$ 差\\
\midrule
case33bw & 33  & 8 & $0.9131$ & 4 & $0.9131$ & $2\times10^{-9}$\\
case69   & 69  & 8 & $0.9092$ & 5 & $0.9092$ & $4\times10^{-10}$\\
case85   & 85  & 9 & $0.8739$ & 5 & $0.8739$ & $5\times10^{-10}$\\
case141  & 141 & 7 & $0.9279$ & 4 & $0.9279$ & $1\times10^{-10}$\\
\bottomrule
\end{tabular}
\end{table}

\subsection{与实现的对应}\label{sec:pf-th-cls-implmap}
\begin{center}\footnotesize
\begin{longtable}{@{}L{6.8cm}L{8.6cm}@{}}
\toprule
理论条目 & 实现状态\\
\midrule
\endfirsthead
\toprule
理论条目 & 实现状态\\
\midrule
\endhead
Y-bus Gauss--Seidel 不动点迭代（定义~\ref{sec:pf-th-cls-def-gs}） &
\implnone\\
Z-bus 阻抗阵法 &
\implnone\\
FDPF 主迭代 $B'\Delta\theta=\Delta P/V$、$B''\Delta V=\Delta Q/V$
（式~\eqref{sec:pf-th-cls-eq:fdpf}） &
\implfull{src/power\_flow/fdpf\_solver.cpp:FDPFSolver::solve}（失配装配见同文件
compute\_power\_injections、compute\_specified\_injections）\\
$B'/B''$ 的 XB 口径常量阵装配（定义~\ref{sec:pf-th-cls-def-xbbx}） &
\implfull{src/power\_flow/admittance\_builder.cpp:build\_susceptance\_matrix}（构型标志
\srcpath{include/hacdcpf/power\_flow/assembly/admittance\_builder.hpp:make\_bp\_flags}、\srcpath{make\_bpp\_flags}）\\
FDPF 常量阵一次分解 + 线性收敛结构
（定理~\ref{sec:pf-th-cls-thm-fdpf}） &
\implpart{Eigen SparseLU 一次性分解（FDPFSolver::solve）；无全局化与 Q 限值切换，
收敛域外的行为无保障}\\
FDPF BX 变体 &
\implnone\\
直流潮流 $B\theta=P$ 降阶求解
（式~\eqref{sec:pf-th-cls-eq:dcpf}，定理~\ref{sec:pf-th-cls-thm-dc}） &
\implfull{src/power\_flow/ac\_linearized\_pf.cpp:solve\_ac\_linearized\_dc}\\
直流潮流移相变压器注入等效 &
\implfull{src/power\_flow/ac\_linearized\_pf.cpp:solve\_ac\_linearized\_dc}（同函数内
$\pm b_k\varphi_k$ 注入装配）\\
直流损耗二次近似与损耗修正重解（式~\eqref{sec:pf-th-cls-eq:loss}） &
\implnone\\
DistFlow 支路潮流方程组
（式~\eqref{sec:pf-th-cls-eq:dfp}--\eqref{sec:pf-th-cls-eq:dfl}） &
\implpart{无独立方程组求解器；辐射支路关系由精确 AC 模型的 BFS 扫描实现，
收敛点与 DistFlow 等价}\\
LinDistFlow 线性化模型 &
\implpart{潮流路径内无；MIP 建模层
src/power\_models/lindistflow\_builder.cpp:solve\_lindistflow（属 power\_models 模块）}\\
前推回代扫描映射（定义~\ref{sec:pf-th-cls-def-bfs}） &
\implfull{src/power\_flow/distribution\_power\_flow.cpp:run\_bfs\_sweep}（BFS 树
build\_bfs\_tree，支路关系 branch\_series\_current、branch\_child\_voltage\_from\_parent，
ZIP 电压依赖 voltage\_dependent\_net\_demand）\\
辐射/单连通前提的 fail-close 检查 &
\implfull{src/power\_flow/distribution\_power\_flow.cpp:solve\_distribution\_pf\_ac\_snapshot}\\
BFS 压缩收敛先验检验（定理~\ref{sec:pf-th-cls-thm-bfs} 的 $L<1$） &
\implpart{无先验条件检验；以 $\max|\Delta V|$ 收敛判据与拓扑 fail-close 兜底}\\
高 $R/X$ 网络的自动方法切换 &
\implnone（方法由用户显式选择，无基于 $R/X$ 的自动分派）\\
参与因子三口径生成（定义~\ref{sec:pf-th-cls-def-dslack}） &
\implfull{src/power\_flow/distributed\_slack\_solver.cpp:create\_participation\_factors}（含
sanitize\_slack\_cfg 校验与 normalize\_factors 归一）\\
事后分摊与带限值重解两口径
（注记~\ref{sec:pf-th-cls-rem-dslack}） &
\implfull{src/power\_flow/distributed\_slack\_solver.cpp:solve\_simplified、solve\_full\_jacobian}（限值重分摊
distribute\_with\_limits，失配复算 compute\_total\_participating\_mismatch\_pu）\\
方程内精确分布式松弛（增广未知量 $\pi$） &
\implnone\\
孤岛分解与每岛参考/slack 自选
（引理~\ref{sec:pf-th-cls-lem-lap} 的算法化） &
\implfull{src/power\_flow/island\_detector.cpp:detect\_islands、extract\_island\_subsystem；自选摇摆
src/power\_flow/adaptive\_solver.cpp:auto\_select\_swing\_bus}\\
死岛零电压哨兵（注记~\ref{sec:pf-th-cls-rem-island}） &
\implfull{src/power\_flow/adaptive\_solver.cpp:set\_dead\_island\_zero}（调度入口
AdaptiveSolver::solve）\\
PTDF/ISF 与 LODF（定义~\ref{def:pf-th-cls-ptdf}、定理~\ref{thm:pf-th-cls-lodf}） &
\implpart{仅市场 N-1 安全校核实现：src/market/market\_simulation.cpp:build\_lodf\_model
构造 PTDF 与 $\psi_{lm}=\rho_{lm}/(1-\rho_{mm})$（lodf\_value）；power\_flow 直流求解器不产出因子}\\
方程内精确分布式松弛加边系统（命题~\ref{prop:pf-th-cls-augslack}，
式~\eqref{sec:pf-th-cls-eq:augslack}） &
\implnone（实现为事后分摊 solve\_simplified 与重复完整重解 solve\_full\_jacobian，见上）\\
FDPF/直流/BFS 收敛性数值基准（表~\ref{tab:pf-th-cls-fdpf}、\ref{tab:pf-th-cls-dc}、
\ref{tab:pf-th-cls-bfs}） &
\implfull{tools/pf\_doc\_benchmark.cpp}（模式 fdpf/dc/bfs）\\
\bottomrule
\end{longtable}
\end{center}

\subsection{参考文献}
{\renewcommand{\section}[2]{}%
\begin{thebibliography}{9}\small
\bibitem{cls:wardhale} J.~B. Ward and H.~W. Hale, Digital computer solution of
power-flow problems, \emph{AIEE Transactions (Power Apparatus and Systems)},
vol. 75, 1956.
\bibitem{cls:stagg} G.~W. Stagg and A.~H. El-Abiad, \emph{Computer Methods in
Power System Analysis}, McGraw-Hill, 1968.
\bibitem{cls:tinney} W.~F. Tinney and C.~E. Hart, Power flow solution by
Newton's method, \emph{IEEE Transactions on Power Apparatus and Systems},
vol. PAS-86, no. 11, 1967.
\bibitem{cls:stott} B.~Stott and O.~Alsa\c{c}, Fast decoupled load flow,
\emph{IEEE Transactions on Power Apparatus and Systems}, vol. PAS-93, no. 3, 1974.
\bibitem{cls:stottreview} B.~Stott, Review of load-flow calculation methods,
\emph{Proceedings of the IEEE}, vol. 62, no. 7, 1974.
\bibitem{cls:vanamerongen} R.~A.~M. van Amerongen, A general-purpose version of
the fast decoupled load flow, \emph{IEEE Transactions on Power Systems},
vol. 4, no. 2, 1989.
\bibitem{cls:monticelli} A.~Monticelli, A.~Garcia and O.~R. Saavedra, Fast
decoupled load flow: hypothesis, derivations, and testing, \emph{IEEE
Transactions on Power Systems}, vol. 5, no. 4, 1990.
\bibitem{cls:ortega} J.~M. Ortega and W.~C. Rheinboldt, \emph{Iterative Solution
of Nonlinear Equations in Several Variables}, Academic Press, 1970.
\bibitem{cls:arrillaga} J.~Arrillaga and N.~R. Watson, \emph{Computer Modelling
of Electrical Power Systems}, 2nd ed., Wiley, 2001.
\bibitem{cls:baranwu} M.~E. Baran and F.~F. Wu, Network reconfiguration in
distribution systems for loss reduction and load balancing, \emph{IEEE
Transactions on Power Delivery}, vol. 4, no. 2, 1989.
\bibitem{cls:shirmohammadi} D.~Shirmohammadi, H.~W. Hong, A.~Semlyen and
G.~X. Luo, A compensation-based power flow method for weakly meshed distribution
and transmission networks, \emph{IEEE Transactions on Power Systems},
vol. 3, no. 2, 1988.
\bibitem{cls:chiangbaran} H.-D. Chiang and M.~E. Baran, On the existence and
uniqueness of load flow solution for radial distribution power networks,
\emph{IEEE Transactions on Circuits and Systems}, vol. 37, no. 3, 1990.
\bibitem{cls:grainger} J.~J. Grainger and W.~D. Stevenson, \emph{Power System
Analysis}, McGraw-Hill, 1994.
\bibitem{cls:kundur} P.~Kundur, \emph{Power System Stability and Control},
McGraw-Hill, 1994.
\bibitem{cls:woodwollenberg} A.~J. Wood, B.~F. Wollenberg and G.~B. Shebl\'e,
\emph{Power Generation, Operation, and Control}, 3rd ed., Wiley, 2014.
\end{thebibliography}}
